% Part: first-order-logic
% Chapter: tableaux
% Section: derivations

\documentclass[../../../include/open-logic-section]{subfiles}

\begin{document}

\iftag{FOL}
      {\olfileid{fol}{tab}{der}}
      {\olfileid{pl}{tab}{der}}

\olsection{\usetoken{P}{tableau}}

\begin{explain}
گفته‌ایم فرض چیست و قواعد استنتاج را به دست داده‌ایم. !!^{tableau}s از
این‌ها به‌طور استقرایی تولید می‌شوند: هر !!{tableau} یا شاخه‌ای یگانه و
متشکل از یک یا چند فرض است، یا از اعمال یکی از قواعد استنتاج بر شاخه‌ای
از !!a{tableau} حاصل می‌شود.
\end{explain}

\begin{defn}[!!^{tableau}]
!!^a{tableau} برای فرض‌های $\sFmla{S_1}{!A_1}$، \dots،
$\sFmla{S_n}{!A_n}$ (که هر $S_i$ یا $\True$ است یا~$\False$)، درختی
متناهی از !!{signed formula}s است که شرایط زیر را برآورده می‌کند:
\begin{enumerate}
\item $n$~!!{signed formula}s بالاییِ درخت،
  $\sFmla{S_i}{!A_i}$ هستند که یکی زیر دیگری قرار می‌گیرند.
\item هر !!{signed formula} در درخت که یکی از فرض‌ها نیست، از کاربرد
  درست یک قاعدهٔ استنتاج بر !!a{signed formula} در شاخهٔ بالای آن حاصل
  می‌شود.
\end{enumerate}
شاخه‌ای از !!a{tableau} \emph{بسته} است هرگاه و تنها هرگاه هم
$\sFmla{\True}{!A}$ و هم~$\sFmla{\False}{!A}$ را دربر داشته باشد؛ وگرنه
\emph{باز} است. !!^a{tableau} که هر شاخه‌اش بسته باشد، \emph{!!{tableau}
بسته} (برای مجموعهٔ فرض‌هایش) است. اگر !!a{tableau} بسته نباشد، یعنی اگر
دست‌کم یک شاخهٔ باز داشته باشد، \emph{باز} است.
\end{defn}

\begin{ex}
  هر مجموعهٔ فرض‌ها به‌تنهایی !!a{tableau} است، اما معمولاً بسته نخواهد
  بود. (بدیهی است که تنها هنگامی بسته است که فرض‌ها از پیش جفتی از
  !!{signed formula}s، یعنی $\sFmla{\True}{!A}$ و
  ~$\sFmla{\False}{!A}$، را دربر داشته باشند.)

  از !!a{tableau} (باز یا بسته) می‌توانیم با اعمال یکی از قواعد استنتاج
  بر !!a{signed formula}~$\sFmla{S}{!A}$ در آن، تابلویی تازه و بزرگ‌تر
  به دست آوریم. قاعده یک یا چند !!{signed formula}s را به انتهای هر شاخه‌ای
  می‌افزاید که رخداد~$\sFmla{S}{!A}$ مورد اعمال قاعده را دربر داشته باشد.

  برای نمونه، فرض $\sFmla{\True}{!A \land \lnot
    !A}$ را در نظر بگیرید. در زیر، !!{tableau} (باز) متشکل فقط از همان
  فرض آمده است:
  \begin{center}
    \begin{tableau}{}
      [\sFmla{\True}{\formula{A} \land \lnot \formula{A}}, just=\TAss]
    \end{tableau}{}
  \end{center}
  با اعمال قاعدهٔ $\TRule{\True}{\land}$ بر فرض، !!{tableau} تازه‌ای از
  آن به دست می‌آوریم. آن قاعده اجازه می‌دهد دو سطر تازه، یعنی
  $\sFmla{\True}{!A}$ و $\sFmla{\True}{\lnot !A}$، را به !!{tableau}
  بیفزاییم:
  \begin{center}
    \begin{tableau}{}
      [\sFmla{\True}{\formula{A} \land \lnot \formula{A}}, just=\TAss
        [\sFmla{\True}{\formula{A}}, just={\TRule{\True}{\land}[1]},
          [\sFmla{\True}{\lnot \formula{A}}, just={\TRule{\True}{\land}[1]}
          ]
        ]
      ]
    \end{tableau}{}
  \end{center}
  هنگام نوشتن !!{tableau}s، قواعد اعمال‌شده را در سمت راست ثبت می‌کنیم
  (برای نمونه، $\TRule{\True}{\land} 1$ یعنی !!{signed formula} آن سطر از
  اعمال قاعدهٔ $\TRule{\True}{\land}$ بر !!{signed formula} سطر~$1$ حاصل
  شده است). این !!{tableau} تازه اکنون !!{signed formula}s دیگری دارد،
  اما قاعده‌ای را فقط بر یکی از آن‌ها ($\sFmla{\True}{\lnot !A}$) می‌توان
  اعمال کرد (در این مورد، قاعدهٔ $\TRule{\True}{\lnot}$). حاصل،
  !!{tableau} بستهٔ زیر است:
  \begin{center}
    \begin{tableau}{}
      [\sFmla{\True}{\formula{A} \land \lnot \formula{A}}, just=\TAss
        [\sFmla{\True}{\formula{A}}, just={\TRule{\True}{\land}[1]}
          [\sFmla{\True}{\lnot \formula{A}}, just={\TRule{\True}{\land}[1]}
            [\sFmla{\False}{\formula{A}}, just={\TRule{\True}{\lnot}[3]}, close]
          ]
        ]
      ]
    \end{tableau}{}
  \end{center}
\end{ex}

\end{document}
